
\subsection{安培力}

\begin{tabularx}{\columnwidth}{XX}
  磁感应强度&$B=\dfrac{F}{Il}$\\
  安培力&$F=BIl\sin\theta$\\
  通电导线受力方向&左手定则\\
  直导线磁场方向&安培定则\\
  螺线管磁场方向&右手螺旋定则\\
\end{tabularx}

\begin{formulanotebox}
$F=BIl$ 只适用于电流方向与磁场方向垂直的情况；一般情形应使用 $F=BIl\sin\theta$，或取垂直于磁场的有效长度。

磁感线是人为建立的模型，磁场是真实存在的物质形态。小磁针 N 极受力方向规定为该点磁场方向。
\end{formulanotebox}

\begin{keypointbox}[title={\strut\textcolor{white}{\bfseries 重点知识点：磁场方向}}]
\textbf{地磁场}：地磁场水平方向大致由地理南方指向地理北方；北半球有向下分量，南半球有向上分量。\\

\textbf{电流磁效应}：直导线周围磁感线为同心圆；螺线管内部磁场类似条形磁铁，内部磁感线从 S 极指向 N 极。\\

\textbf{安培力方向}：左手定则中四指指向电流方向，磁感线穿过掌心，大拇指指向安培力方向。
\end{keypointbox}

\begin{casetypebox}[title={\strut\textcolor{white}{\bfseries 常见题型：通电导线受力与平衡}}]
\begin{itemize}[leftmargin=1.5em]
  \item 先判断磁场方向，再用左手定则判断安培力方向。
  \item 同向平行电流相互吸引，反向平行电流相互排斥。
  \item 导线在磁场中平衡时，安培力常与重力、拉力、弹力共同列平衡方程。
  \item 若导线与磁场不垂直，只取垂直磁场方向的电流长度。
\end{itemize}
\end{casetypebox}
